% ======================================
% 06-environments.tex - 定理/定义/例题环境定义
% ======================================
% 2024-04-14 更新：为每种环境类型定义独立计数器

% ========== 独立计数器定义 ==========
% 定理类环境共享 theorem 计数器
\newcounter{theorem}[chapter]
\renewcommand{\thetheorem}{\thechapter.\arabic{theorem}}

% 定义类环境使用 definition 计数器
\newcounter{definition}[chapter]
\renewcommand{\thedefinition}{\thechapter.\arabic{definition}}

% 例题类环境使用 example 计数器
\newcounter{example}[chapter]
\renewcommand{\theexample}{\thechapter.\arabic{example}}

% 练习类环境使用 exercise 计数器
\newcounter{exercise}[chapter]
\renewcommand{\theexercise}{\thechapter.\arabic{exercise}}

% 真题环境使用 realexam 计数器
\newcounter{realexam}[chapter]
\renewcommand{\therealexam}{\thechapter.\arabic{realexam}}

% ========== 定理类环境 ==========
\newtcolorbox{thmBox}[3][]{
  commonbox,
  colback=LogicBg,
  colbacktitle=LogicColor,
  before title={\refstepcounter{theorem}},
  title={#2 \thetheorem\ifstrempty{#3}{}{（#3）}},
  #1
}

\newenvironment{theorem}[1][]{\begin{thmBox}{定理}{#1}}{\end{thmBox}}
\newenvironment{lemma}[1][]{\begin{thmBox}{引理}{#1}}{\end{thmBox}}
\newenvironment{proposition}[1][]{\begin{thmBox}{命题}{#1}}{\end{thmBox}}
\newenvironment{corollary}[1][]{\begin{thmBox}{推论}{#1}}{\end{thmBox}}
\newenvironment{property}[1][]{\begin{thmBox}{性质}{#1}}{\end{thmBox}}

% ========== 定义类环境 ==========
\newtcolorbox{defBox}[3][]{
    commonbox,
    colback=DefBg,
    colbacktitle=DefColor,
    before title={\refstepcounter{definition}},
    title={#2 \thedefinition\ifstrempty{#3}{}{（#3）}},
    #1
}

\newenvironment{definition}[1][]{\begin{defBox}{定义}{#1}}{\end{defBox}}
\newenvironment{axiom}[1][]{\begin{defBox}{公理}{#1}}{\end{defBox}}

% ========== 例题类环境 ==========
\newtcolorbox{exBox}[3][]{
    commonbox,
    colback=ExBg,
    colbacktitle=ExColor,
    before title={\refstepcounter{example}},
    title={#2 \theexample\ifstrempty{#3}{}{（#3）}},
    #1
}

\newtcolorbox{longexBox}[3][]{
    breakable,
    commonbox,
    colback=ExBg,
    colbacktitle=ExColor,
    before title={\refstepcounter{example}},
    title={#2 \theexample\ifstrempty{#3}{}{（#3）}},
    #1
}

\newenvironment{example}[1][]{\begin{exBox}{例题}{#1}}{\end{exBox}}
\newenvironment{longexample}[1][]{\begin{longexBox}{例题}{#1}}{\end{longexBox}}

% ========== 练习环境（独立计数器）==========
\newtcolorbox{exerciseBox}[3][]{
    commonbox,
    colback=ExBg,
    colbacktitle=ExColor,
    before title={\refstepcounter{exercise}},
    title={#2 \theexercise\ifstrempty{#3}{}{（#3）}},
    #1
}

\newtcolorbox{longexerciseBox}[3][]{
    breakable,
    commonbox,
    colback=ExBg,
    colbacktitle=ExColor,
    before title={\refstepcounter{exercise}},
    title={#2 \theexercise\ifstrempty{#3}{}{（#3）}},
    #1
}

\newenvironment{exercise}[1][]{\begin{exerciseBox}{练习}{#1}}{\end{exerciseBox}}
\newenvironment{longexercise}[1][]{\begin{longexerciseBox}{练习}{#1}}{\end{longexerciseBox}}

% ========== 真题环境 ==========
\newtcolorbox{realexamBox}[2][]{
    commonbox,
    colback=ExBg,
    colbacktitle=ExColor,
    before title={\refstepcounter{realexam}},
    title={真题\ \therealexam\ifstrempty{#2}{}{（#2）}},
    #1
}

\newtcolorbox{longrealexamBox}[2][]{
    breakable,
    commonbox,
    colback=ExBg,
    colbacktitle=ExColor,
    before title={\refstepcounter{realexam}},
    title={真题\ \therealexam\ifstrempty{#2}{}{（#2）}},
    #1
}

\newenvironment{realexam}[1][]
{\begin{realexamBox}{#1}}
{\end{realexamBox}}

\newenvironment{longrealexam}[1][]
{\begin{longrealexamBox}{#1}}
{\end{longrealexamBox}}

% ========== 注与证明 ==========
\newenvironment{note}{
    \par\vspace{0.5em}\noindent
    \textbf{\color{NoteColor} 注：}
}{
    \par\vspace{0.5em}
}

\renewenvironment{proof}{
    \par\vspace{0.5em}\noindent
    \textbf{\color{LogicColor} 证明：}
}{
    \hfill \textcolor{LogicColor}{\rule{0.6em}{0.6em}} \par
}

% ======================================
% 交互元素环境 - 新增
% ======================================

% ========== 学生思路环境 ==========
\newtcolorbox{thinkingBox}[1][]{
    enhanced,
    colback=ThinkingBg,
    colframe=ThinkingColor,
    coltitle=white,
    colbacktitle=ThinkingColor,
    fonttitle=\bfseries,
    attach boxed title to top left={xshift=5mm, yshift*=-\tcboxedtitleheight/2},
    boxed title style={
        frame hidden,
        rounded corners,
        arc=3pt,
        top=1mm, bottom=1mm, left=2mm, right=2mm
    },
    boxrule=1pt,
    sharp corners=downhill, rounded corners, arc=8pt,
    left=5mm, right=5mm, top=4mm, bottom=4mm,
    before skip=1.2em, after skip=1em,
    title={\faUser\ 学生思路\ifstrempty{#1}{}{：#1}},
    fontupper=\kaishu\linespread{1.5}\selectfont
}

\newenvironment{studentthinking}[1][]
{\begin{thinkingBox}{#1}}
{\end{thinkingBox}}

% ========== 教师点睛环境 ==========
\newtcolorbox{insightBox}[1][]{
    enhanced,
    colback=InsightBg,
    colframe=InsightColor,
    coltitle=white,
    colbacktitle=InsightColor,
    fonttitle=\bfseries,
    attach boxed title to top right={xshift=-5mm, yshift*=-\tcboxedtitleheight/2},
    boxed title style={
        frame hidden,
        rounded corners,
        arc=3pt,
        top=1mm, bottom=1mm, left=2mm, right=2mm
    },
    boxrule=1.5pt,
    sharp corners=downhill, rounded corners, arc=8pt,
    left=5mm, right=5mm, top=4mm, bottom=4mm,
    before skip=1.2em, after skip=1em,
    title={\faLightbulbO\ 教师点睛\ifstrempty{#1}{}{：#1}},
    fontupper=\bfseries\color{InsightColor}\linespread{1.5}\selectfont
}

\newenvironment{teacherinsight}[1][]
{\begin{insightBox}{#1}}
{\end{insightBox}}

% ========== 变式训练环境 ==========
\newcounter{variation}[chapter]
\renewcommand{\thevariation}{\thechapter.\arabic{variation}}

\newtcolorbox{variationBox}[2][]{
    enhanced,
    colback=VariationBg,
    colframe=VariationColor,
    coltitle=white,
    colbacktitle=VariationColor,
    fonttitle=\bfseries,
    attach boxed title to top left={xshift=5mm, yshift*=-\tcboxedtitleheight/2},
    boxed title style={
        frame hidden,
        rounded corners,
        arc=3pt,
        top=1mm, bottom=1mm, left=2mm, right=2mm
    },
    boxrule=1pt,
    sharp corners=downhill, rounded corners, arc=8pt,
    left=5mm, right=5mm, top=4mm, bottom=4mm,
    before skip=1.2em, after skip=1em,
    before title={\refstepcounter{variation}},
    title={\faCodeFork\ 变式训练 \thevariation\ifstrempty{#2}{}{（#2）}},
    fontupper=\linespread{1.5}\selectfont,
    #1
}

\newenvironment{variation}[1][]
{\begin{variationBox}{#1}}
{\end{variationBox}}

% ========== 提示环境（轻量级） ==========
\newtcolorbox{hintBox}{
    enhanced,
    colback=HintBg,
    colframe=HintColor,
    boxrule=0.5pt,
    arc=4pt,
    left=3mm, right=3mm, top=2mm, bottom=2mm,
    before skip=0.8em, after skip=0.8em,
    fontupper=\small\kaishu\color{HintColor}
}

\newenvironment{hint}{
\begin{hintBox}
\textbf{\color{HintColor}\faHandORight\ 提示：}
}
{\end{hintBox}}

% ========== 常见错误环境 ==========
\newtcolorbox{errorBox}[1][]{
    enhanced,
    colback=CommonErrorBg,
    colframe=CommonErrorColor,
    coltitle=white,
    colbacktitle=CommonErrorColor,
    fonttitle=\bfseries,
    attach boxed title to top left={xshift=5mm, yshift*=-\tcboxedtitleheight/2},
    boxed title style={
        frame hidden,
        rounded corners,
        arc=3pt,
        top=1mm, bottom=1mm, left=2mm, right=2mm
    },
    boxrule=1pt,
    sharp corners=downhill, rounded corners, arc=8pt,
    left=5mm, right=5mm, top=4mm, bottom=4mm,
    before skip=1.2em, after skip=1em,
    title={\faExclamationTriangle\ 常见错误\ifstrempty{#1}{}{：#1}},
    fontupper=\linespread{1.5}\selectfont
}

\newenvironment{commonerror}[1][]
{\begin{errorBox}{#1}}
{\end{errorBox}}

% ========== 步骤拆解环境 ==========
\newenvironment{steps}[1][步骤详解]
{
    \par\vspace{0.5em}
    \noindent\textbf{\color{ExColor}#1：}
    \begin{enumerate}[label=\textbf{\color{ExColor}\arabic*.}]
}
{
    \end{enumerate}
    \par\vspace{0.5em}
}

% ========== 知识卡片环境 ==========
\newtcolorbox{knowledgecard}[1][]{
    enhanced,
    colback=white,
    colframe=DefColor,
    boxrule=1pt,
    arc=6pt,
    left=4mm, right=4mm, top=3mm, bottom=3mm,
    before skip=1em, after skip=1em,
    fontupper=\small,
    title={\faBookmark\ 知识卡片\ifstrempty{#1}{}{：#1}},
    coltitle=DefColor,
    fonttitle=\bfseries\small,
    attach title to upper={\par\vspace{1mm}}
}

\newenvironment{card}[1][]
{\begin{knowledgecard}{#1}}
{\end{knowledgecard}}
